\cdcsection{Maquettes}

Les maquettes suivantes donnent une première idée de l’interface finale. Elles sont volontairement filaires. Elles fixent l’organisation des écrans et les informations affichées, pas leur charte graphique. Les valeurs « x » remplacent des données réelles.

\cdcmaquettes{assets/maquettes/m1-recherche}{Recherche de boutiques (client)}%
             {assets/maquettes/m2-enregistrement}{Enregistrement d’un appareil (réparateur)}

Le client décrit son problème et voit sur la carte les boutiques capables de le traiter (F4, F5). Le réparateur enregistre l’appareil et photographie son état initial depuis son smartphone (F2).

\cdcmaquettes{assets/maquettes/m3-validation}{Validation de l’état initial (client)}%
             {assets/maquettes/m4-suivi}{Suivi et paiement (client)}

Le client valide ou conteste l’état initial avant tout diagnostic (F2). Il suit ensuite l’avancement de sa réparation, consulte le prix et le délai confirmés, puis choisit son moyen de paiement (F3, F6).

\cdcmaquette{assets/maquettes/m5-tableau-de-bord}{Tableau de bord du propriétaire (ordinateur)}

Le propriétaire suit l’activité de sa boutique, ses revenus, la répartition des paiements et les performances de ses réparateurs (F7, F12).
